\subsection{Research Gap}
\label{sec:relwork:gap}
The two halves surveyed above do not meet, and each opening below is narrower than a
claim that nobody has coarsened a circuit graph.
% Original: The two halves surveyed above do not meet. The openings are specific, and each is narrower
% than a claim that nobody has coarsened a circuit graph.

\begin{itemize}
    \item \textbf{Unseen designs are seldom the tested condition.} Prediction on designs
    never seen in training is the field's stated motivation, yet only 4 of 63 surveyed
    papers report matched seen and unseen numbers (\ref{sec:relwork:ml4eda:protocols}): under a
    random or recipe-disjoint split, every test design also appears in training under
    another recipe, so an unmeasured share of the reported score is design recognition,
    not structural inference. The standard split trains on 16 small designs and tests on
    8 large ones \citep{chowdhury2021openabc}, confounding novelty with size, and a model
    motivated on the gap need not publish the seen-design baseline \citep{deng2024hoga}.

    \item \textbf{Circuit models already reduce their input and never measure it.} The
    representation-learning line trains on extracted subcircuits of tens to a few
    thousand gates \citep{li2022deepgate,deepgate4_2025}, a reduction applied at corpus
    construction and neither named nor evaluated as one
    (\ref{sec:relwork:ml4eda:representation}). What it costs is unknown.

    \item \textbf{Graph reduction has not been evaluated on AIGs.} The taxonomy is
    developed and surveyed \citep{hashemi2024survey,shabani2023summarization}, but no
    method in it has been measured against a logic-synthesis target. Transfer cannot be
    assumed: a generic method cuts or merges under an undirected criterion, and a single
    severed edge removes gates from the fanin cone~\eqref{eq:prelim:cone} of everything
    downstream (\ref{sec:prelim:aig:distinct}).

    \item \textbf{Coarsening results in this domain carry no compression-matched
    control.} CTS-Bench reports memory and speed gains for one bespoke heuristic and
    compares no second method \citep{cts_bench2026}, so a gain cannot be attributed to
    the merge criterion rather than the compression.

    \item \textbf{The two equivalence literatures are not connected.} Structural
    hashing, FRAIG and maximal fanout-free cone decomposition are coarsenings by merge
    map \citep{kuehlmann2002robust,mishchenko2005fraig,mishchenko2006rewriting}. The
    exactness framework deciding which an encoder can absorb without loss
    \citep{bollen2023exact} has never been turned on them.
\end{itemize}
% Original bullets:
% Unseen designs: Prediction on designs the model has never seen is the field's stated motivation.
% Matched seen and unseen-design numbers appear in 4 of the 63 papers surveyed here
% (\ref{sec:relwork:ml4eda:protocols}). Under a random or recipe-disjoint split every test
% design also appears in training under another recipe, so an unmeasured share of the
% reported score is design recognition rather than structural inference. The standard
% design-disjoint split on the reference corpus trains on 16 small designs and tests on 8
% large ones \citep{chowdhury2021openabc}, which confounds novelty with size, and a model
% motivated on the gap need not publish the seen-design number its degradation would be
% measured against \citep{deng2024hoga}.
%
% Circuit models: The representation-learning line trains on extracted subcircuits of tens
% to a few thousand gates \citep{li2022deepgate,deepgate4_2025}. Extraction is a reduction,
% applied at corpus construction, and it is neither named as one nor evaluated as one
% (\ref{sec:relwork:ml4eda:representation}). What it costs is unknown.
%
% Graph reduction on AIGs: The taxonomy is developed and surveyed
% \citep{hashemi2024survey,shabani2023summarization}, and no method in it has been measured
% against a logic-synthesis target. Transfer cannot be assumed, for a mechanical reason: a
% generic method cuts or merges under an undirected criterion, and a single severed edge
% removes gates from the fanin cone~\eqref{eq:prelim:cone} of everything downstream of it
% (\ref{sec:prelim:aig:distinct}).
%
% Compression-matched control: CTS-Bench reports memory and speed gains for one bespoke
% heuristic and compares no second method \citep{cts_bench2026}. Absent a control at the
% same ratio, a gain cannot be attributed to the merge criterion rather than to the
% compression.
%
% Equivalence literatures: Structural hashing, FRAIG and maximal fanout-free cone
% decomposition are coarsenings by merge map
% \citep{kuehlmann2002robust,mishchenko2005fraig,mishchenko2006rewriting}, and the
% exactness framework that decides which of them an encoder can absorb without loss
% \citep{bollen2023exact} has never been turned on them.

After CTS-Bench, being first to coarsen in EDA cannot be claimed unqualified, so the claim
made here is narrower. This thesis is the first to (i) benchmark multiple coarsening
families, exact refinement-based, GNN-aware, and domain-specific, against a
compression-matched random control for GNN training in EDA, (ii) do so on AIGs and
optimizability regression, (iii) bring provably lossless compression into an EDA GNN
setting, (iv) connect AIG-native equivalence reduction to the generic summarization
framework, and (v) run a controlled three-protocol split comparison with per-design error
distributions. Spectral and hashing-based coarsening are surveyed
in~\ref{sec:relwork:reduction:summarization} but not run. Results are compared against
SynthNet \citep{chowdhury2021openabc}, HOGA \citep{deng2024hoga} and DeepGate4
\citep{deepgate4_2025} on the prediction task, and against CTS-Bench
\citep{cts_bench2026} on the memory, speed and accuracy trade-off.
% Original: After CTS-Bench, being first to coarsen in EDA cannot be claimed unqualified, and the claim
% made here is narrower. This thesis is the first to (i) benchmark multiple coarsening
% families, exact refinement-based, GNN-aware, and domain-specific, against a
% compression-matched random control for GNN training in EDA, (ii) do so on AIGs and
% optimizability regression, (iii) bring provably lossless compression into an EDA GNN
% setting, (iv) connect AIG-native equivalence reduction to the generic summarization
% framework, and (v) run a controlled three-protocol split comparison with per-design error
% distributions on this task. Spectral and hashing-based coarsening are surveyed
% in~\ref{sec:relwork:reduction:summarization} but not run, and are no part of the claim. The
% eventual results are compared against SynthNet \citep{chowdhury2021openabc}, HOGA
% \citep{deng2024hoga} and DeepGate4 \citep{deepgate4_2025} on the prediction task, and
% against CTS-Bench \citep{cts_bench2026} on the shape of the memory, speed and accuracy
% trade-off.
